\lyseis{Γ1.1}{Απογραφή, πληθυσμός, κατηγορικά και ποσοτικά δεδομένα}
\sol{1} α) Κατηγορικό. β) Ποσοτικό διακριτό. γ) Κατηγορικό. δ) Ποσοτικό συνεχές. ε) Ποσοτικό συνεχές. στ) Ποσοτικό διακριτό. ζ) Κατηγορικό.
\sol{2} Πληθυσμός: οι $120$ μαθητές της Α΄ τάξης. Δείγμα: τα $20$ παιδιά που ρωτήθηκαν. Μεταβλητή: ο αριθμός των βιβλίων που διάβασε κάθε παιδί. Τα δεδομένα είναι ποσοτικά και διακριτά, γιατί ο αριθμός των βιβλίων είναι ακέραιος.
\sol{3} α) Απογραφή (εξετάζονται και τα $25$ παιδιά). β) Δείγμα ($500$ από $12\,000$). γ) Δείγμα (ένα σταφύλι από όλον τον αμπελώνα). δ) Απογραφή (μετρούνται όλα τα παιδιά της τάξης).
\sol{4} Το σύνολο είναι $20$ και οι σχετικές συχνότητες βγαίνουν: τοστ $\frac{7}{20}=35\%$, φρούτο $\frac{5}{20}=25\%$, κουλούρι $\frac{4}{20}=20\%$, γιαούρτι $\frac{2}{20}=10\%$, τίποτα $\frac{2}{20}=10\%$. Άθροισμα: $35+25+20+10+10=100\%$.
\sol{5} α) Ο πιο συνηθισμένος τρόπος είναι με τα πόδια ($9$ παιδιά) και ο πιο σπάνιος με ποδήλατο ($3$ παιδιά). β) $25-8=17$ παιδιά. γ) $\frac{3}{25}=0{,}12=12\%$.
\sol{6} Οι συχνότητες είναι $12$, $8$, $6$, $4$, $2$ και το σύνολο $12+8+6+4+2=32$ ζώα. Στο ραβδόγραμμα οι ράβδοι έχουν αυτά τα ύψη, ίδιο πλάτος και κενά μεταξύ τους (ελεύθερη κατασκευή).
\sol{7} α) Λάθος (μπορεί να είναι και αντικείμενα ή ζώα). β) Σωστό. γ) Σωστό. δ) Λάθος (είναι κατηγορικό). ε) Σωστό.
\sol{8} Ο τρόπος β). Με την κλήρωση κάθε παιδί του σχολείου έχει την ίδια πιθανότητα να ρωτηθεί και το δείγμα είναι αντιπροσωπευτικό. Στο α) οι φίλοι του Στέλιου μοιάζουν στα γούστα και στο γ) δεν ρωτιούνται όσοι δεν πηγαίνουν στο κυλικείο εκείνο το διάλειμμα.
\sol{9} Λευκό $8$ ($40\%$), γκρι $5$ ($25\%$), μαύρο $4$ ($20\%$), κόκκινο $2$ ($10\%$), άλλο $1$ ($5\%$). Σύνολο $20$ και $100\%$. Το ραβδόγραμμα έχει ύψη $8$, $5$, $4$, $2$, $1$ (ελεύθερη κατασκευή).
\sol{10} α) Η τιμή «$1$ βιβλίο» με συχνότητα $9$. β) $8+6+3=17$ παιδιά. γ) $\frac{9}{30}=30\%$. δ) $0\cdot4+1\cdot9+2\cdot8+3\cdot6+4\cdot3=0+9+16+18+12=55$ βιβλία.
\sol{11} Χρώμα: κατηγορικά. Μνήμη (GB): ποσοτικά διακριτά. Βάρος (g): ποσοτικά συνεχή. Αδιάβροχο: κατηγορικά.
\sol{12} Το σκάκι έχει $40-(10+8+6+12)=40-36=4$. Οι σχετικές συχνότητες: μουσική $25\%$, χορός $20\%$, ζωγραφική $15\%$, ρομποτική $30\%$, σκάκι $10\%$. Άθροισμα $100\%$.
\sol{13} Όχι αυτόματα. Το $32\%$ ισχύει για τα $25$ παιδιά της Α$3$. Για να ισχύει για όλο το σχολείο, το δείγμα πρέπει να είναι αντιπροσωπευτικό, δηλαδή να έχει παιδιά από όλες τις τάξεις, διαλεγμένα με κλήρωση. Μια μόνο τάξη μπορεί να έχει διαφορετικές προτιμήσεις.
\sol{14} Ελεύθερη διατύπωση. Παράδειγμα: πληθυσμός οι μαθητές της Α΄ Γυμνασίου της πόλης, δείγμα $40$ παιδιά που επιλέγονται με κλήρωση από διάφορα σχολεία, μεταβλητή οι ώρες ύπνου, δεδομένα ποσοτικά συνεχή.
\sol{15} Ο Ηλίας έχει δίκιο. Το άθροισμα είναι $40+35+30=105\%$, ενώ πρέπει να είναι $100\%$. Άρα υπάρχει λάθος στον υπολογισμό ή στην καταμέτρηση.
\sol{16} Ελεύθερη απάντηση. Παραδείγματα: α) «Ποιο είναι το αγαπημένο σου χρώμα;» β) «Πόσα αδέλφια έχεις;» γ) «Πόσο ζυγίζει η τσάντα σου;»
\sol{17} Ελεύθερη κατασκευή. Το άθροισμα των συχνοτήτων πρέπει να βγαίνει $10$ και των σχετικών συχνοτήτων $100\%$.
\sol{18} Ελεύθερη απάντηση. Ο μαθητής πρέπει να ονομάσει σωστά τον πληθυσμό (ποιους αφορά το διάγραμμα), τη μεταβλητή και το είδος των δεδομένων.
